\documentclass[11pt]{article}
\usepackage{amsmath,amssymb}
\usepackage{graphicx}
\usepackage[hidelinks]{hyperref}

\newcommand{\R}{\mathbb{R}}
\DeclareMathOperator{\tr}{tr}

\title{A Small Example}
\author{Ada Lovelace}
\date{\today}

\begin{document}
\maketitle

\begin{abstract}
This document shows what the \LaTeX{} editor draws in its Visual mode: headings, \textbf{bold}, \emph{emphasis},
\texttt{code}, math, lists, references and a figure.
\end{abstract}

\section{Introduction}
\label{sec:intro}

Let $f\colon \R^n \to \R$ be a smooth function. Its gradient $\nabla f$ points
where $f$ grows fastest --- a fact we use in Section~\ref{sec:method}.
Euler's identity, $e^{i\pi} + 1 = 0$, links five constants.

% A comment stays in the source but is dimmed in Visual mode.

\section{Method}
\label{sec:method}

For a square matrix $A$ we have
\begin{equation}
  \tr(AB) = \tr(BA), \qquad \det(AB) = \det A \cdot \det B.
  \label{eq:trace}
\end{equation}
Equation~\eqref{eq:trace} holds for any $A, B \in \R^{n\times n}$.

\subsection{Steps}

\begin{enumerate}
  \item Pick a starting point $x_0$.
  \item Repeat $x_{k+1} = x_k - \eta \nabla f(x_k)$ until it settles.
  \item Report $x_k$.
\end{enumerate}

The step size matters:
\begin{itemize}
  \item too small, and it crawls;
  \item too large, and it diverges\footnote{See any textbook on optimisation.}.
\end{itemize}

\section{Conclusion}

Gradient descent is simple and often enough~\cite{knuth1984}:
\[
  \min_{x \in \R^n} f(x) = \sum_{i=1}^{m} \bigl(y_i - g(x)\bigr)^2 .
\]

\begin{thebibliography}{9}
\bibitem{knuth1984} D.~E. Knuth, \emph{The \TeX book}, Addison-Wesley, 1984.
\end{thebibliography}

\end{document}
